\sect{Математическое описание задачи}{background}
%
В настоящем разделе формализуется задача отслеживания партитуры,
последовательно выводятся алгоритмы Dynamic Time Warping и On-line
Time Warping, а также описывается скрытая марковская модель, лежащая в
основе нашего трекера. Далее приведены рекуррентные формулы Forward и
Viterbi и кратко обсуждается схема комбинации HMM и OLTW.

Партитура задаётся последовательностью
$S=\{(p_i,\,t_i^{\text{on}},\,t_i^{\text{off}})\}_{i=1}^{N}$
из $N$ нот с высотой $p_i\in\{0,\dots,127\}$ (стандарт MIDI) и
номинальными моментами начала и конца звучания. Исполнение поступает
каузально как поток
$P=\{(q_j,\,\tau_j^{\text{on}},\,\tau_j^{\text{off}})\}_{j=1}^{M}$.
Задача отслеживания состоит в построении отображения~\eqref{eq:alignment},
введённого в~\Secref{introduction}, по информации только о прошлых
наблюдениях $q_1,\dots,q_j$.

Определим локальную функцию стоимости как абсолютную разность высот
$d(s_i,q_j) = |p_i - q_j|$ в полутонах. В большинстве случаев такая
мера достаточна; ошибки октавного типа отдельно обрабатываются на
этапе очистки входных данных (см.~\Secref{realtime}).

\textit{Dynamic Time Warping}~\cite{SakoeChiba1978,Mueller2007} строит
матрицу накопленных стоимостей $D\in\Real^{N\times M}$ согласно
рекуррентной формуле
\begin{equation}
D(i,j) = d(s_i,q_j) + \min\bigl\{D(i-1,j),\,D(i,j-1),\,
                                  D(i-1,j-1)\bigr\},
\eqlab{dtw-recursion}
\end{equation}
с граничным условием $D(1,1)=d(s_1,q_1)$. Оптимальное выравнивание
получается обратным ходом из ячейки $(N,M)$. Полоса Сакоэ-Чиба
ограничивает допустимые ячейки условием $|i-j|\le R$ и снижает
сложность до $O(NR)$, что при $R\approx 0{,}1\,N$ соответствует
ускорению на порядок при сохранении точности на ровном темпе.

\par\addvspace{2pt}
{\centering
\refstepcounter{figure}\figlab{dtw-matrix}%
\resizebox{0.85\columnwidth}{!}{% =====================================================================
% TikZ figure: DTW cost matrix with Sakoe-Chiba band and optimal path
% Used in: sections/03-background.tex
% =====================================================================
\begin{tikzpicture}[x=6mm,y=6mm,font=\scriptsize]

  % grid 6x6
  \foreach \i in {0,...,6}
    \draw[black!50] (0,\i) -- (6,\i);
  \foreach \j in {0,...,6}
    \draw[black!50] (\j,0) -- (\j,6);

  % Sakoe-Chiba band (|i-j| <= 2)
  \fill[black!10]
    (0,0) -- (3,0) -- (6,3) -- (6,6) -- (3,6) -- (0,3) -- cycle;
  \foreach \i in {0,...,6}
    \draw[black!50] (0,\i) -- (6,\i);
  \foreach \j in {0,...,6}
    \draw[black!50] (\j,0) -- (\j,6);

  % cost matrix values (illustrative)
  \foreach \i/\j/\v in {%
    0/0/0, 0/1/3, 0/2/5, 0/3/7,
    1/0/2, 1/1/0, 1/2/2, 1/3/4, 1/4/6,
    2/0/4, 2/1/2, 2/2/0, 2/3/2, 2/4/4, 2/5/6,
    3/1/4, 3/2/2, 3/3/0, 3/4/2, 3/5/4, 3/6/6,
    4/2/4, 4/3/2, 4/4/0, 4/5/2, 4/6/4,
    5/3/4, 5/4/2, 5/5/0, 5/6/2,
    6/4/4, 6/5/2, 6/6/0
  } \node at (\j+0.5,\i+0.5) {\v};

  % optimal path
  \draw[ultra thick]
    (0.5,0.5) -- (1.5,1.5) -- (2.5,2.5) -- (3.5,3.5)
    -- (4.5,4.5) -- (5.5,5.5);
  \foreach \x/\y in {0.5/0.5,1.5/1.5,2.5/2.5,3.5/3.5,4.5/4.5,5.5/5.5}
    \fill (\x,\y) circle (1.2pt);

  % axes
  \draw[->,thick] (0,-0.2) -- (6.4,-0.2)
    node[below right] {Партитура $i$};
  \draw[->,thick] (-0.2,0) -- (-0.2,6.4)
    node[above left] {Исполнение $j$};

  % corner labels
  \node[below=2pt] at (0.5,-0.2) {1};
  \node[below=2pt] at (5.5,-0.2) {$N$};
  \node[left=2pt]  at (-0.2,0.5) {1};
  \node[left=2pt]  at (-0.2,5.5) {$M$};
\end{tikzpicture}
}\par
\vspace{2pt}
\footnotesize
\figurename~\thefigure. Матрица накопленных стоимостей DTW для
пятинотного фрагмента из прелюдии Шопена. Серым выделена полоса
Сакоэ-Чиба, чёрной ломаной~--- оптимальный путь выравнивания.\par}
\addvspace{4pt}

В каузальной постановке полная матрица $D$ недоступна. Метод On-line
Time Warping~\cite{Dixon2005} обрабатывает столбец за столбцом по мере
поступления новой ноты $q_j$ и поддерживает текущую оценку
выравнивания $\hat\varphi(j)$. Алгоритм управляется счётчиком
\textit{RunCount}: если выбор минимума в~\Eqref{dtw-recursion} приводит
к увеличению строки $i$ более $R_{\max}$ раз подряд, обработка
принудительно сдвигается по столбцам, что предотвращает «застревание»
на одной позиции партитуры. Псевдокод OLTW и подробный анализ
сложности приведены в~\Secref{appendix-oltw}.

Скрытая марковская модель ставит в соответствие каждой ноте партитуры
скрытое состояние $s_i\in\{1,\dots,N\}$ и описывает наблюдения
$o_t$ (либо высоту MIDI-ноты, либо вектор признаков аудио). Модель
характеризуется тройкой $\lambda=(A,B,\pi)$~\cite{Rabiner1989}: матрица
переходов $A=(A_{ij})$, эмиссионные распределения $b_i(o)$ и начальное
распределение $\pi_i$. Для score-following естественна полосная
структура матрицы переходов:
\begin{equation}
A_{ij} = \begin{cases}
p_{\text{stay}},    & j = i,\\
p_{\text{advance}}, & j = i+1,\\
p_{\text{skip}},    & j = i+2,\\
0,                  & \text{иначе},
\end{cases}
\eqlab{transition-matrix}
\end{equation}
с условием нормировки
$p_{\text{stay}}+p_{\text{advance}}+p_{\text{skip}}=1$. Самозамыкание
описывает ферматы и удлинённые ноты, переход $i\!\to\!i+2$~--- редкий
пропуск ноты исполнителем. Эмиссия для MIDI-входа моделируется
гауссовой плотностью
$b_i(o)=\mathcal{N}(o;\,p_i,\sigma^2)$, для аудиовхода~--- свёрточным
модулем (см.~\Secref{cnn-corner}).

\par\addvspace{2pt}
{\centering
\refstepcounter{figure}\figlab{hmm-states}%
\resizebox{\columnwidth}{!}{% =====================================================================
% TikZ figure: HMM state diagram for score following
% Used in: sections/03-background.tex (subsec:hmm)
% =====================================================================
\begin{tikzpicture}[node distance=12mm and 14mm]

  \foreach \i/\note [count=\xi] in {1/C4, 2/E4, 3/G4, 4/A4, 5/C5} {
    \node[state] (s\i) at (\xi*16mm, 0) {$s_{\i}$};
    \node[font=\tiny, below=1mm of s\i] {\note};
  }

  % advance arrows (top)
  \foreach \a/\b in {1/2, 2/3, 3/4, 4/5}
    \draw[arrow] (s\a) -- node[above, font=\tiny]
                          {$0.6$} (s\b);

  % stay self-loops
  \foreach \i in {1,...,5}
    \path (s\i) edge[arrow, loop above, looseness=8,
                     in=110, out=70]
                node[font=\tiny] {$0.3$} (s\i);

  % skip arrows (dashed, bend below)
  \foreach \a/\c in {1/3, 2/4, 3/5}
    \draw[arrow, dashed] (s\a) to[bend right=35]
         node[below, font=\tiny] {$0.1$} (s\c);

  % observations
  \foreach \i [count=\xi] in {1,2,3,4,5} {
    \node[obs, below=14mm of s\i] (o\i) {$o_{\i}$};
    \draw[arrow, faded] (s\i) -- (o\i);
  }

  % legend
  \node[font=\scriptsize, anchor=north west,
        text width=42mm] at ($(s5.east)+(8mm,5mm)$) {%
    $\to$\ $p_{\text{advance}}\!=\!0.6$\\
    $\circlearrowleft$\ $p_{\text{stay}}\!=\!0.3$\\
    $\dashleftarrow$\ $p_{\text{skip}}\!=\!0.1$};

\end{tikzpicture}
}\par
\vspace{2pt}
\footnotesize
\figurename~\thefigure. Диаграмма пятисостояний HMM для фрагмента
$S=(\text{C4, E4, G4, A4, C5})$ с типичными значениями
$p_{\text{stay}}=0{,}3$, $p_{\text{advance}}=0{,}6$,
$p_{\text{skip}}=0{,}1$.\par}
\addvspace{4pt}

Каузальное оценивание положения исполнителя сводится к рекуррентному
вычислению нормированного forward-распределения:
\begin{equation}
\alpha_t(j) \propto b_j(o_t)\sum_{i=1}^{N}\alpha_{t-1}(i)\,A_{ij},
\qquad \hat\varphi(t) = \argmaxop_{j}\alpha_t(j),
\eqlab{forward-recursion}
\end{equation}
с инициализацией $\alpha_0(j)=\pi_j$. Сложность одного шага составляет
$O(N)$ благодаря разреженности матрицы~$A$. Соответствующая Viterbi-
рекурсия и её псевдокод вынесены в~\Secref{appendix-hmm}.

Стандартный вариант HMM моделирует длительность пребывания в одном
состоянии геометрическим распределением
$P(d)=p_{\text{stay}}^{d-1}(1-p_{\text{stay}})$, что некорректно для
партитур с большой разницей длительности нот. В реализации использован
подход с явной длительностью: состояние расширяется до пары $(i,d)$,
где $d$~--- счётчик такта в текущей ноте, а вероятность перехода
$p_{\text{advance}}(d,d_{\text{exp}})$ зависит от ожидаемой
длительности $d_{\text{exp}}$ согласно партитуре.

Комбинация HMM и OLTW использует сильные стороны обеих моделей. HMM
выдаёт калиброванное апостериорное распределение, по которому удобно
оценивать уверенность через
$c_t = \max_i\alpha_t(i)$. Когда $c_t<\theta_{\text{conf}}$
(в наших экспериментах $\theta_{\text{conf}}=0{,}3$), выходная оценка
переключается на OLTW, и трекер запускает процедуру повторной
синхронизации. Такая схема рассматривается далее в~\Secref{realtime}.
